% ---------------------------------------------------------------------------
%  UNIT III -- Chapter 10
% ---------------------------------------------------------------------------
\chapter{Unit III --- Understanding Harmony in Nature and Existence}
\chaptermark{Unit III --- Chapter 10}

\section{A. Nature, Harmony and the Four Orders}

% ------------------------------------------------------------------------ 50
\Q{50}{Define harmony in nature. Why is it important? Explain with
  examples.}{\tg{SEE}\tg{CIE 3 QB}\mmk{5}}

\ans Nature is the combination of all that is in the solid, liquid or gas
state; in other words, the \textbf{aggregate of all the mutually interacting
units --- big or small, sentient or insentient --- together can be called
nature}. These units are infinite in number, and we could easily observe that
there exists a \textbf{dynamic balance, a self-regulation} among all these
units. \textbf{This self-regulation is harmony or balance in nature.}

The law of nature has a unique cause-and-effect system which must be
understood in order to be in harmony with the natural law of things. Natural
harmony is necessary for the following reasons:

\begin{apts}
\item It is necessary to solve the problem of \textbf{global warming}, and the
      depletion of non-renewable natural resources can be avoided.
\item Natural harmony with trees cures many problems --- reduction of wind
      velocity, energy savings, companion planting, development of an
      eco-subsystem in terms of establishing a forest garden, and reduction of
      building heat.
\item It is possible to achieve natural harmony in the establishment,
      maintenance and management of educational institutions like schools,
      colleges and universities.
\item One can understand the depths of harmony and alignment in nature by
      contemplating and reflecting upon the natural order. It is possible to
      unravel the mystery of the natural synthesis in the midst of ongoing
      chaos at the material plane.
\end{apts}

\textbf{How to create it.} First of all we have to understand the importance
of nature for our existence, and then our responsibility towards nature. When
we take from nature, we should also plan to give back to it. Renewable energy
sources such as solar power and wind are perhaps the best methods to begin
moving in this direction. There can be harmony only if each subsystem of the
entire creation can strike a mutually satisfying relationship with every other
subsystem, without disturbing the other's peaceful existence and without
hindering its own growth. To create harmony in nature, we have to work in the
direction of the development of mankind from \textbf{animal consciousness to
human consciousness} --- and this entails working for the \textbf{right
understanding}.

% ------------------------------------------------------------------------ 51
\Q{51}{What are the four orders of nature? Briefly explain
  them.}{\tg{CIE 3}\tg{SEE}\tg{Scheme}\mmk{5}}

\ans All the physical objects in the solid, liquid or gas state, living or
non-living, are collectively termed nature --- the aggregate of all mutually
interacting units, big or small, sentient or insentient. These units are
infinite in number and there exists a dynamic balance and self-regulation
among them. The concept of the four orders provides a framework for
understanding the natural world as an \textbf{interconnected system composed
of distinct yet interdependent categories}:

\begin{apts}
\item \textbf{Material Order (Padartha Avastha \,/\, Bhautik Avastha).} The
      big land masses of the continents, gigantic water bodies like oceans and
      seas, mountains and rivers, the atmosphere above, the heaps of metals
      and minerals below, the dense gases and fossil fuels deep below the
      surface of the earth --- all fall into the material order. Beyond the
      earth it is visible as stars, planets, moons and several astronomical
      bodies. It comprises all \textbf{non-living elements} --- soil, water,
      air, minerals and other inanimate objects --- and forms the foundational
      base for all other forms of life. It follows natural laws like
      gravitation, thermodynamics and chemical reactions; these elements are
      stable and do not possess the ability to grow or reproduce.
\item \textbf{Pranic \,/\, Plant \,/\, Bio Order (Prana Avastha \,/\, Jeevan
      Avastha).} Our land mass is covered with grass and small shrubs which
      form the lining on the entire soil; shrubs, plants and trees form huge
      forests, along with the flora in the ocean. This order includes all
      \textbf{living but non-conscious} entities --- plants, trees and other
      vegetation. It is responsible for sustaining life by producing oxygen,
      absorbing carbon dioxide, and providing food and shelter to animals and
      humans. It is characterised by growth, reproduction and the ability to
      interact with the environment through processes like photosynthesis.
      However, it does not have self-awareness.
\item \textbf{Animal Order (Jeeva Avastha \,/\, Chetan Avastha).} Animals and
      birds form the third largest order. It includes all \textbf{conscious
      living beings} --- animals, birds, insects and other creatures that
      possess sensory abilities. Animals play a crucial role in maintaining
      ecological balance by being part of the food chain, aiding in
      pollination, seed dispersal and other natural processes. The animal
      order is characterised by \textbf{instincts, sensory perception and
      consciousness}. Animals can experience pleasure, pain and emotions, but
      are primarily driven by instincts rather than rational thinking.
\item \textbf{Human \,/\, Knowledge Order (Manav Avastha \,/\, Gyana
      Avastha).} Humans are the smallest order. The human order encompasses
      all human beings, who are \textbf{conscious, self-aware and capable of
      rational thought}. Humans have the unique ability to reflect on their
      actions, make ethical decisions, create culture and innovate. They also
      have the responsibility to nurture and maintain harmony among the other
      orders. The human order is defined by self-awareness, the ability to
      think rationally, creativity, and the pursuit of knowledge.
\end{apts}

\begin{keybox}
\textbf{Note the relative quantity:} the material order is far greater than
the plant\,/\,bio order, which is far greater than the animal order, which is
far greater than the human order.
\end{keybox}

\textbf{Conclusion.} The four orders of nature --- physical, bio, animal and
human --- are interconnected and interdependent. Each order plays a vital role
in maintaining the balance and harmony of the natural world. Understanding
these orders helps us appreciate the complexity of existence and encourages
responsible and sustainable living.

% ------------------------------------------------------------------------ 52
\Q{52}{What are the four orders in nature? Describe their activity,
  innateness, natural characteristic, basic activity and
  conformance.}{\tg{CIE 3 QB}\mmk{8--10}}

\ans The four orders can be distinctly recognised in terms of the following
five parameters:

\renewcommand{\arraystretch}{1.3}
\arrayrulecolor{rulegrey}
\begin{longtable}{P{27mm} P{31mm} P{31mm} P{31mm} P{31mm}}
\rowcolor{slate} & \hdr{Material Order} & \hdr{Pranic \,/\, Bio Order} &
  \hdr{Animal Order} & \hdr{Human Order} \\ \hline \endfirsthead
\rowcolor{slate} & \hdr{Material Order} & \hdr{Pranic \,/\, Bio Order} &
  \hdr{Animal Order} & \hdr{Human Order} \\ \hline \endhead
\textbf{Things} & Soil, air, water, metals, minerals &
  Plants, herbs, trees; animal bodies and human bodies &
  Animal body $+$ `I' & Human body $+$ `I' \\ \hline
\textbf{Activity} \newline (Kriya) & Composition\slash Decomposition &
  Composition\slash Decomposition $+$ Respiration &
  \emph{Body:} Comp.\slash Decomp., Respiration\newline
  \emph{`I':} $+$ Selection (assuming) &
  \emph{Body:} Comp.\slash Decomp., Respiration\newline
  \emph{`I':} Desiring, Thinking, Selecting --- with the possibility of
  Realization \& Understanding \\ \hline
\textbf{Innateness} \newline (Dharana) & Existence & Existence $+$ Growth &
  (Existence $+$ Growth) in body $+$ \textbf{Will to live} in `I' &
  (Existence $+$ Growth) in body $+$ \textbf{Will to live with happiness} in
  `I' \\ \hline
\textbf{Natural Charac\-teristic} \newline (Svabhava) &
  Composition\slash Decomposition &
  Comp.\slash Decomp. $+$ Nurture\slash Worsen &
  (Comp.\slash Decomp., Nurture\slash Worsen) in body $+$
  \textbf{Non-cruelty\slash Cruelty} in `I' &
  (Comp.\slash Decomp., Nurture\slash Worsen) in body $+$
  \textbf{Perseverance, Bravery, Generosity} in `I' \\ \hline
\textbf{Basic Activity} & Recognising, Fulfilment &
  Recognising, Fulfilment &
  \emph{Body:} Recognising, Fulfilment\newline
  \emph{`I':} Assuming, Recognising, Fulfilment &
  \emph{Body:} Recognising, Fulfilment\newline
  \emph{`I':} Knowing, Assuming, Recognising, Fulfilment \\ \hline
\textbf{Conformance} \newline (Anushangita) & Constitution conformance &
  Seed conformance & Breed conformance & Sanskar conformance \\ \hline
\end{longtable}

\begin{note}
\textbf{\sffamily This single table is the highest-yield thing to memorise for
Unit III.} Most Unit III questions are a row or a column of it. Q53 is the
\emph{Innateness} row, Q54 the \emph{Svabhava} row, Q55 the
\emph{Conformance} row, Q56 the \emph{Activity} and \emph{Basic activity}
rows, Q57 the Pranic and Animal columns, Q58 the Animal and Human columns.
\end{note}

% ------------------------------------------------------------------------ 53
\Q{53}{What do you mean by `innateness' (dharana)? What is the innateness in
  the four orders?}{\tg{CIE 3 QB}\mmk{5}}

\ans \textbf{Innateness (dharana)} means the qualities which are innate to the
unit. Each unit in existence exhibits an innateness, an intrinsic quality that
\textbf{cannot be separated from it}. This is intrinsic to the unit.

\begin{abul}
\item \textbf{Material order --- Existence.} When we burn coal and only some
      ash is left and the smoke has gone out, it is not that the fundamental
      particles in the coal have ``ceased to exist''. They may not be visible
      to the eye at that moment, but they continue to exist in the form of
      other matter or gases. We cannot destroy matter; we can only convert it
      from one form to another. Thus \emph{``to exist''} is intrinsic to all
      material --- we cannot separate the `existence' of a thing from the
      thing itself.
\item \textbf{Pranic \,/\, plant order --- Existence $+$ Growth.} Because the
      pranic order is a development of the material order, it also has the
      innateness of `existence'. In addition it exhibits `growth'. This
      principle cannot be separated from any unit of this order: if you have a
      plant, you cannot stop it from growing. The only way to stop it is by
      cutting it --- but when you do that, it ceases to belong to the pranic
      order and, instead, decays and belongs to the material order.
\item \textbf{Animal order --- Existence $+$ Growth $+$ Will to live.} The
      animal body is a development of the pranic order and therefore inherits
      `existence' and `growth' at the level of the body, which is
      physico-chemical in nature. In addition, all units in this order have
      the \textbf{`will to live'} in `I'. No unit in this order can be
      separated from this will to live.
\item \textbf{Human (knowledge) order --- Existence $+$ Growth $+$ Will to
      live with happiness.} In the human being, `existence' and `growth' are
      fundamentally present in the body, just as in the animal body. At the
      level of `I', however, in addition to the `will to live', a human
      being's innateness is the \textbf{`will to live with happiness'}.
\end{abul}

% ------------------------------------------------------------------------ 54
\Q{54}{What is the natural characteristic (svabhava) of a unit? Explain the
  svabhava of the human order and how it helps in living with harmony.}%
  {\tg{CIE 3}\tg{SEE}\tg{SEE 2025}\tg{Scheme}\mmk{5--7}}

\ans When we look at the different orders in nature, we find that each order
has a certain value --- in a fundamental way, this is the
\textbf{`usefulness' or `participation'} of the order in existence. This value
or participation is also referred to as the \textbf{natural characteristic or
svabhava}. The characteristic the order displays is natural to itself.

\begin{cmptable}{Order}{Svabhava}
Material order & Composition \,/\, decomposition \\ \hline
Pranic \,/\, plant order & Composition \,/\, decomposition, and to
  \textbf{nurture or worsen} other pranic units \\ \hline
Animal order & \emph{Body:} composition\,/\,decomposition,
  nurture\,/\,worsen\newline
  \emph{`I':} \textbf{non-cruelty (akrurata) and cruelty (krurata)}. Cruelty
  means the feeling that it can fulfil its needs through violence and
  forcefulness. Cows may largely live with non-cruelty, while tigers and lions
  may exhibit cruelty. \\ \hline
Human order & \emph{Body:} composition\,/\,decomposition,
  nurture\,/\,worsen\newline
  \emph{`I':} \textbf{perseverance (dhirata), bravery (virata) and generosity
  (udarata)} \\ \hline
\end{cmptable}

\textbf{Svabhava of the human order in detail.} As in the case of animals, the
human body also belongs to the plant\,/\,bio order and hence has the same
svabhava as the pranic order --- it either nurtures or worsens other pranic
units. (When I digest a vegetable, I absorb the plant and it worsens, while my
body is nurtured.) The svabhava of the self (`I') in human beings is:

\begin{abul}
\item \textbf{Perseverance (dhirata)} --- being assured that the
      all-encompassing solution is to understand and live in harmony at all
      levels of existence, and living with this commitment \emph{without any
      perturbation}.
\item \textbf{Bravery (virata)} --- being assured of this solution, and being
      \emph{ready to help the other have the right understanding}. This is the
      commitment to help the other understand harmony and to live at all
      levels of existence.
\item \textbf{Generosity (udarata)} --- being assured of this solution, and
      being \emph{ready to invest myself, my body and my wealth} to help the
      other have the right understanding.
\end{abul}

\textbf{How this helps in living with harmony.} Human beings are not living as
per this natural characteristic; even though we have a svabhava, we are not
living according to it. This is the basic reason for the contradiction and
conflict that we see in the human being, and what leads to a state of
unhappiness. Only when we live according to our basic human characteristics as
mentioned above do we have a \textbf{definite character}; otherwise it is not
definite, it is uncertain --- unlike the other three orders, which are always
according to their svabhava.

\begin{note}
\textbf{\sffamily The scheme's alternative wording.} The 2024 CIE~3 scheme
expresses the same idea in general terms: the svabhava of the human order
revolves around \emph{self-awareness and consciousness, the desire for
knowledge and understanding, the ability to make rational decisions, the
aspiration for justice and ethics, and creativity and innovation}. Either
formulation is acceptable; \textbf{the textbook triad (dhirata, virata,
udarata) is the safer answer} --- give it first, and add the scheme's list
only if marks remain.
\end{note}

% ------------------------------------------------------------------------ 55
\Q{55}{What do you mean by `conformance'? Explain the conformance in the four
  orders. \emph{(Also set as:} What is sanskaar? Explain its effects or the
  conformance of the human order.\emph{)}}%
  {\tg{CIE 3 QB}\tg{Impr.\ CIE 2025}\tg{SEE 2025}\mmk{5--7}}

\ans Each unit conforms through the principle of \textbf{conformance
(anushangita)}. It means \textbf{how the continuity of the fundamental nature
of the unit is preserved}.

\begin{abul}
\item \textbf{Material order --- Constitution conformance.} The continuity of
      the fundamental nature of a material unit is preserved through physical
      and chemical processes. Take iron: each atom of iron conforms to the
      constitutional structure of `iron'. There is no atom of iron that will
      be unlike another atom of iron --- if it were, we would not call it
      iron. The same holds for oxygen, nitrogen, gold, silver, aluminium: all
      conform to the constitution of their kind.
\item \textbf{Plant \,/\, bio order --- Seed conformance.} A neem seed will
      always sprout a neem plant. Its fruits, its leaves, the taste and colour
      of the leaves --- all this basic information of every neem plant is
      stored in the seed. Thus the plant is always as the seed: \emph{``as the
      seed, thus the plant''}. This seed-conformance method is the mechanism
      by which the continuity of a plant species is maintained in
      nature\,/\,existence.
\item \textbf{Animal order --- Breed conformance.} A cow is always like a cow,
      and a dog is always like a dog. Animals conform to their lineage --- how
      animals are, and their behaviour, is according to the lineage they come
      from. This breed-conformance method is the mechanism by which the
      continuity of an animal species is maintained in nature\,/\,existence.
\item \textbf{Human (knowledge) order --- Sanskar conformance.} We humans are
      \emph{not} according to our lineage or race, as animals are. We may pick
      up something from our parents as we grow up, but we are usually very
      different from them in many ways. We humans are according to our
      \textbf{imagination} --- according to our desires, thoughts and
      selections in `I'. These can come from past memories, our parents, the
      environment, the media --- anywhere. In the case of humans we can say
      \emph{``as the education, so the human''}. Together we call these
      \textbf{sanskara}; hence a human being conforms to his or her sanskar.
\end{abul}

\begin{keybox}
\textbf{\sffamily The `what is sanskaar?' version.} \emph{Sanskaar} is the sum
of the desires, thoughts and selections in `I' --- our imagination --- from
whatever source it has been picked up. \textbf{Its effect} is that it, and not
lineage, is what the human being conforms to. This is why the human is the
only order whose conduct is \emph{not definite}: an animal is always according
to its breed, but a human is according to sanskaar, which can be right or
wrong. It is also why education matters so much in this course --- ``as the
education, so the human'' --- and why right understanding, once ensured, fixes
conduct at its source.
\end{keybox}

% ------------------------------------------------------------------------ 56
\Q{56}{How is the activity in the human order different from that in the
  animal and plant order? Explain the activity in the four
  orders.}{\tg{CIE 3 QB}\mmk{5}}

\ans An activity means something that `has motion' and\,/\,or `has a result'.
We can understand this in two ways: things that are \emph{visibly moving} (a
spinning top, a moving bus) are active; and things that are \emph{visibly
stationary} are also active --- e.g. the activity of a chair is that the wood
is interacting with the environment and, as a result, decays with time. All
units around us, including ourselves, are active all the time; in the activity
there is a state or configuration and motion simultaneously.

\begin{abul}
\item \textbf{Material order.} All material things can be understood as an
      activity of units coming together to form a bigger unit ---
      \emph{composition} (e.g. a chair made of smaller pieces of wood). Bigger
      units can also separate to form smaller units --- \emph{decomposition}
      (a wooden chair decays after a few years). Thus any unit in the material
      order is an activity of composition\,/\,decomposition.
\item \textbf{Plant \,/\, bio order.} Composition\,/\,decomposition $+$
      \textbf{respiration}. Not only do plants compose (forming new plants)
      and decompose (decaying), they are also breathing or pulsating, which we
      call respiration.
\item \textbf{Animal order.} Two aspects:
      \emph{Body --- physico-chemical activities:} the body displays the same
      activities as the plant --- composition\,/\,decomposition and
      respiration; hence we say the body belongs to the plant\,/\,bio order.
      \emph{`I' --- conscious activities:} `I' is a unit that has the ability
      of \textbf{assuming}. Animals make assumptions. If we have a dog and a
      stranger comes into the house, the dog may start barking at him. If this
      person stays at our house, the dog may stop barking at him but will
      continue to bark at other strangers. The dog's \emph{assumption} about
      this person has changed, due to which the way it responds has changed.
      This consciousness or faculty of assuming is not in the body; the body
      belongs to the plant\,/\,bio order and just responds to physico-chemical
      inputs.
\item \textbf{Human (knowledge) order.} The activities in the human body are
      similar to those in the animal body --- composition\,/\,decomposition
      and respiration. When it comes to consciousness or `I', however, the
      human displays more than just an ability to select or make choices as
      animals do. In human beings, `I' has the activities of \textbf{desiring,
      thinking, and selecting\,/\,tasting, with a possibility or need for
      realization and understanding}. \textbf{Only humans have this need to
      know}, and that is why it is called \emph{gyana avastha --- the
      knowledge order}.
\end{abul}

\textbf{Basic activity.} In the material and pranic orders there is only
\textbf{recognising and fulfilment}; such units do not have the activities of
assuming and knowing. Hydrogen and oxygen recognise the relation to each other
and combine to form water. A brick and another brick have a definite relation,
recognise it, and get arranged to form a building. A plant recognises the
relation with sun and water and fulfils it by acting accordingly. There is no
selection involved --- a plant does not choose to turn or not to turn to the
sun. When we look at animals and humans, we find \emph{selection} taking
place; in animals there is additionally \emph{assuming}, and in humans there
is additionally \emph{knowing}.

% ------------------------------------------------------------------------ 57
\Q{57}{Explain the difference and similarities between the Pranic order and
  the Animal order. What is the relation between the two
  orders?}{\tg{CIE 3}\tg{Scheme}\mmk{5}}

\ans
\begin{cmptable}{Pranic Order (Bio Order)}{Animal Order}
All forms of plant life, including trees, shrubs, grasses and other
  vegetation. &
All animal species, including mammals, birds, reptiles, fish, insects and
  others. \\ \hline
\textbf{Lacks conscious awareness} and sensory perception; operates primarily
  on biological and chemical processes. &
\textbf{Possesses conscious awareness}, sensory perception and instincts;
  animals can perceive their environment, feel pain, and respond to stimuli.
  \\ \hline
Engages in growth, reproduction, photosynthesis and other life-sustaining
  processes without the need for consciousness or decision-making. &
Engages in \textbf{instinct-driven behaviours} such as hunting, mating and
  social interaction, guided by sensory input and biological needs. \\ \hline
Relies on sunlight, water and nutrients from the soil to produce energy
  through photosynthesis. &
Relies on consuming plants (herbivores) or other animals (carnivores) for
  energy. \\ \hline
Typically stationary, with movements limited to growth or response to
  environmental stimuli (e.g.\ plants growing towards light). &
Typically mobile, with the ability to move from one place to another in search
  of food, shelter or mates. \\ \hline
\end{cmptable}

\textbf{Similarities:}
\begin{apts}
\item \textbf{Life processes:} both orders are living entities that undergo
      biological processes such as growth, reproduction and respiration.
\item \textbf{Interdependence:} both are interconnected and interdependent
      within ecosystems. Plants provide oxygen and food for animals, while
      animals contribute to plant reproduction through pollination, seed
      dispersal and nutrient cycling.
\item \textbf{Ecological roles:} both play crucial roles in maintaining
      ecological balance. Plants are primary producers, converting sunlight
      into energy, while animals often act as consumers, maintaining the
      balance of populations within ecosystems.
\end{apts}

\textbf{Relation between the two orders:}
\begin{apts}
\item The \textbf{animal body is a development of the pranic order} --- the
      body of the animal belongs to the plant\,/\,bio order and hence has the
      same `usefulness' or value (nurture\,/\,worsen) as the pranic order.
      Consequently, the animal order inherits the innateness of the pranic
      order --- \textbf{existence and growth} --- and, in addition, has the
      \textbf{will to live} in `I'.
\item The Pranic Order can therefore be seen as the broader concept
      encompassing the life force that animates living bodies, while the
      Animal Order is the \textbf{co-existence of that pranic body with the
      self (`I')}. In this way the Pranic Order provides a more holistic view
      that includes the scientific categorisation found in the Animal Order.
\end{apts}

% ------------------------------------------------------------------------ 58
\Q{58}{Explain the differences and similarities between the animal order and
  the human order. What is the relation between the two?}{\tg{CIE 3 QB}\mmk{5}}

\ans Both orders are the \textbf{co-existence of a body (pranic order) and the
self (`I' $=$ consciousness)}. The animal order is the co-existence of the
animal body and `I'; the human order is the co-existence of the human body and
`I'. That is the fundamental similarity, and the body-level activities
(composition\,/\,decomposition and respiration) are identical in both. The
differences are entirely at the level of `I':

\renewcommand{\arraystretch}{1.3}
\arrayrulecolor{rulegrey}
\begin{longtable}{P{28mm} P{60mm} P{60mm}}
\rowcolor{slate}\hdr{Parameter} & \hdr{Animal Order} & \hdr{Human Order}
  \\ \hline \endfirsthead
\rowcolor{slate}\hdr{Parameter} & \hdr{Animal Order} & \hdr{Human Order}
  \\ \hline \endhead
\textbf{Activity in `I'} & \textbf{Assuming} --- the dog that stops barking at
  a person who stays in the house, but keeps barking at other strangers, has
  changed its assumption. &
  \textbf{Desiring, thinking, selecting\,/\,tasting}, with a possibility or
  need for \textbf{realization and understanding}. Only humans have this need
  to know --- hence \emph{gyana avastha}. \\ \hline
\textbf{Innateness} & Existence $+$ growth (body) $+$ \textbf{will to live}
  (`I') &
  Existence $+$ growth (body) $+$ \textbf{will to live with happiness} (`I')
  \\ \hline
\textbf{Svabhava} & Nurture\,/\,worsen (body) $+$ \textbf{non-cruelty\,/\,
  cruelty} (`I') &
  Nurture\,/\,worsen (body) $+$ \textbf{perseverance, bravery, generosity}
  (`I') \\ \hline
\textbf{Basic activity} & Recognising, fulfilment (body); assuming,
  recognising, fulfilment (`I') &
  Recognising, fulfilment (body); \textbf{knowing}, assuming, recognising,
  fulfilment (`I') \\ \hline
\textbf{Conformance} & \textbf{Breed conformance} --- a cow is always like a
  cow, a dog always like a dog; behaviour is according to lineage. &
  \textbf{Sanskar conformance} --- we are according to our desires, thoughts
  and selections, which can come from memories, parents, environment or media.
  \emph{``As the education, so the human.''} \\ \hline
\textbf{Definiteness of conduct} & \textbf{Definite} --- the animal is always
  according to its breed. &
  \textbf{Not definite today} --- because we are not living according to our
  svabhava; this is the basic reason for contradiction and conflict in the
  human being. \\ \hline
\end{longtable}

\textbf{Relation.} The human order is a further development of the animal
order: the human body belongs to the plant\,/\,bio order just as the animal
body does, and the human inherits existence, growth and the will to live. What
is added in the human order is the activity of \textbf{knowing}, and
correspondingly the innateness of the \textbf{will to live with happiness}.
This addition is precisely what makes the human being \emph{responsible} for
maintaining harmony among the other three orders.

\section{B. Interconnectedness, Recyclability and Self-organisation}

% ------------------------------------------------------------------------ 59
\Q{59}{How will you show interconnectedness and mutual fulfilment in the four
  orders of nature with examples? How is the human order responsible to the
  other three orders?}{\tg{CIE 3 QB}\tg{Impr.\ CIE 2026}\mmk{5}}

\ans In nature, all the units are connected to each other and fulfilling each
other. A human being is related to all other human beings, and is connected to
all the material units in existence, becoming aware of it as he\,/\,she starts
exploring.

\begin{abul}
\item \textbf{Material order and Plant\,/\,Bio order.} The material order
      provides the nutrients to the plant\,/\,bio order in the form of soil,
      minerals, etc., while the plant\,/\,bio order decays and forms more
      nutrients, thus enriching the soil. The plant\,/\,bio order also decays
      into substances like oil and coal, stored deep within the earth as
      protection against the heat from the molten core and from the sun.
      Plants help move the nutrients through the various layers of the soil;
      the roots hold the soil together and prevent erosion; plants produce
      oxygen\,/\,carbon dioxide and thus help in the movement of the material
      order. There is mutual interdependency and co-existence here.
\item \textbf{Material, Plant\,/\,Bio and Animal orders.} The material order
      provides the basis for the movement of all animals, birds and fishes.
      Water, oxygen and other gases are necessities for both plants and
      animals. At the same time, the animal order helps enrich the soil with
      its excreta, which gives nutrients to plants. The plant\,/\,bio order
      provides food for animals, birds and fishes; the animal order helps in
      the pollination of the flowers of the pranic order.
\item \textbf{All four orders including the Human order.} We humans also have
      a natural acceptance to be mutually fulfilling to these three orders.
      \textbf{However, we are not able to ensure this mutual fulfilment.} We
      are dependent on the material order for soil, minerals and metals, but
      only end up polluting the soil and depleting the fossil fuels. We are
      dependent on plants for our food and for holding together the larger
      ecosystem, but we have destroyed forests and multiple species of plants
      and herbs. We are dependent on animals to carry out our production and
      transportation activities, but have made many species extinct, and are
      today known for our cruelty towards animals.
\end{abul}

\textbf{Conclusion.} We can see that there is interconnectedness and mutual
fulfilment in all the orders of nature \textbf{except the human order. We have
to work on this.} The purpose of science and technology is to facilitate the
cyclic processes in nature and make human beings more and more fulfilling to
the other entities. If we understand the processes in nature, we can design
our production systems through the application of science and technology in
such a way that this mutual fulfilment is better ensured, rather than
disturbed.

% ------------------------------------------------------------------------ 60
\Q{60}{Explain how there is recyclability and self-regulation in nature, with
  examples. How is it useful for sustainable production?}%
  {\tg{CIE 3 QB}\tg{Impr.\ CIE 2025}\mmk{5--6}}

\ans There are several \textbf{cyclical processes} that we can see in nature.
For example, the cycle of water --- evaporating, condensing and precipitating
back to water, giving the weather phenomena. The cycles keep these materials
self-regulated on the earth.

\textbf{Examples of self-regulation:}

\begin{abul}
\item \textbf{Forest:} the growth of trees takes place in a way such that the
      amount of soil, plants and animals remains conserved. It never happens
      that the number of trees shoots up and there is a lack of soil for the
      trees. The appropriateness of the conditions for the growth of both
      plants and animals is self-regulated in nature, keeping the population
      proportions naturally maintained.
\item \textbf{Breeds:} in a single breed of animals, the number of males and
      females generated through procreation is such that the continuity of the
      species is ensured by itself. (This happens with humans too, but inhuman
      practices have led to disproportionate numbers of men and women.)
\end{abul}

These two characteristics --- \textbf{cyclical nature and self-regulation} ---
provide us with clues of the harmony that is in nature.

\textbf{Usefulness for sustainable production.} Since nature's systems are
cyclic and mutually enriching, our production systems must be designed within
that framework --- i.e. the \textbf{avartansheel (cyclical) process}, which is
cyclic and ensures that every unit is enriched. Our processes today are acyclic
(open-ended) and deplete resources: when we burn coal we can never again
produce it, and burning it in enormous quantities pollutes the atmosphere.
Recognising recyclability and self-regulation tells us \emph{how} to produce
--- extend the production system already inherently present in nature, so that
resources are continuously renewed and replenished, and production becomes
sustainable rather than depleting.

% ------------------------------------------------------------------------ 61
\Q{61}{How can we say that `nature is self-organised and in space
  self-organisation is available'?}{\tg{CIE 3}\tg{SEE}\tg{Scheme}\mmk{4--5}}

\ans \textbf{Every unit is an organisation.} A unit recognises other units and
combines to form a bigger organisation. Starting from the atom to the big
galaxy, this organisation goes on, as a self-organisation. At every level we
get a self-organisation: sub-atomic particles recognise each other and come
together to form atoms; cells recognise each other and form organisations like
organs and a body; planetary bodies, solar systems and galaxies are still
bigger organisations. \textbf{We are not organising it. We are not supplying
it organisation from outside.}

\textbf{In the human being.} We are self-organised at the level of the body
--- we are not doing anything for the coordination between the heart, kidneys,
lungs, eyes, brain, hands and legs; all these function together. Our input is
needed only to provide the required nutrition and to assist the body when we
fall sick or get injured. At the level of `I', however, we are \emph{not}
self-organised --- but, being in space, \textbf{self-organisation is
available} to the self. That is why we are in pursuit of happiness, which is
essentially being in harmony; whenever we are not in harmony, we are unhappy.

All the units of the four orders are self-organised. No one is organising them
from outside; no one is supplying this organisation. This self-organisation is
available to units \textbf{being in space}. Hence, for space, we say
\textbf{``self-organisation is available''}.

\textbf{Supporting points (from the CIE 3 scheme):}

\begin{abul}
\item \textbf{Natural balance:} ecosystems regulate themselves through food
      chains, nutrient cycles and population-control mechanisms, ensuring that
      no single species dominates to the detriment of others.
\item \textbf{Autonomous processes:} the water cycle, the carbon cycle and the
      growth of plants occur without external intervention, governed by
      natural laws.
\item \textbf{Feedback mechanisms:} predator--prey relationships keep animal
      populations in check; plants adjust their growth based on the
      availability of sunlight and nutrients.
\item \textbf{Role of space:} space is not just a void but an all-encompassing
      entity that provides the medium for everything to exist and interact. It
      allows for the free movement of energy, matter and information, and
      provides the context within which natural laws operate.
\item \textbf{Resilience and adaptation:} ecosystems recover from disturbances
      such as fires or floods through natural regeneration, guided by the
      inherent organisation within nature.
\end{abul}

\section{C. Existence as Co-existence}

% ------------------------------------------------------------------------ 62
\Q{62}{Define existence. ``Existence is co-existence of mutually interacting
  units in all-pervasive space.'' Explain.}{\tg{CIE 3}\tg{SEE}\tg{Scheme}%
  \mmk{5}}

\ans All the units together constitute nature. All the units of nature exist
in space, which is an important reality to understand. \textbf{Existence is
nothing but nature submerged in space.}

\begin{keybox}
\textbf{Existence $=$ Space $+$ Units (in space)}

\vspace{2pt}
Since nature consists of the four orders: \textbf{Existence $=$ Nature
submerged in space}, where Nature $=$ the four orders (Material,
Plant\,/\,Bio or Pranic, Animal and Human Order).
\end{keybox}

A unit is something that is limited in size --- from a small blade of hair to
the biggest planets, all are limited in size, i.e. bounded on six sides. We
can recognise them; they are countable. But there is another reality called
space. We normally don't pay attention to it because it is not a unit --- we
can't touch it or smell it, we just see through it. But the fact that we can't
touch it does not mean it does not exist. Space exists everywhere. When we
look at existence around us, the first thing we see is space; then we see the
units in space, and between every two units there is space.

\textbf{Explaining the statement (as per the CIE 3 scheme).} Existence is
fundamentally the co-existence of various units that interact with one another
within the infinite expanse of space. Each unit --- whether it is a physical
entity like a rock, a biological entity like a plant, or a conscious being
like a human --- is \textbf{not isolated}; rather, it exists in relation to
and interacts with other units.

\begin{apts}
\item \textbf{Co-existence:} all entities or units, be they material or
      living, do not exist in isolation. They inherently exist together,
      forming a cohesive whole. This concept highlights the interconnectedness
      of all elements in the universe.
\item \textbf{Mutual interaction:} units in existence are in constant
      interaction with one another. These interactions are governed by natural
      laws that ensure balance and harmony within the system. For example,
      plants (bio order) interact with the physical environment (soil, water)
      and animals, facilitating the flow of energy and matter.
\item \textbf{All-pervasive space:} space is not merely an empty void but the
      foundation in which all units exist and interact. It is all-pervasive,
      meaning it is present everywhere, allowing these interactions to occur
      seamlessly.
\item \textbf{Holistic perspective:} the idea of co-existence emphasises a
      holistic view of the universe, where the well-being of each unit is tied
      to the well-being of the whole. Disruptions in this balance can lead to
      disharmony, affecting the entire system.
\item \textbf{Sustainability:} recognising existence as co-existence leads to
      an understanding of the need for sustainable living. When humans
      acknowledge their role as part of this interconnected system, they are
      more likely to engage in practices that promote balance and harmony.
\end{apts}

\textbf{In conclusion}, existence is a dynamic process of co-existence where
all units, through their interactions within the vast expanse of space,
contribute to the overall harmony of the universe. This perspective encourages
us to live in a way that respects and nurtures these connections.
\textbf{Existence is in the form of co-existence. It is in harmony. We don't
have to create this harmony --- it already exists. We only have to understand
it to be in it.}

% ------------------------------------------------------------------------ 63
\Q{63}{What do you mean by co-existence? \emph{(Also set as:} Give any two
  definitions for the term co-existence.\emph{)}}%
  {\tg{SEE}\tg{CIE 2/3 QB}\tg{SEE 2025}\tg{Impr.\ CIE 2026}\mmk{1--6}}

\ans \textbf{Co-existence in nature means there is a relationship and
complementarity among all the entities in nature, including human beings.}
Co-existence is a state in which two or more groups are living together while
respecting their differences and resolving their conflicts non-violently. It
has been defined in numerous ways:

\begin{apts}
\item To exist together (in time or space) and to exist in mutual tolerance.
\item To learn to recognise and live with difference.
\item To have a relationship between persons or groups in which none of the
      parties is trying to destroy the other.
\item To exist together (in time or place) and to exist in mutual tolerance.
\end{apts}

The world is full of diversity --- there are different nations, cultures,
religions, communities, languages and beliefs. The beauty of existence can only
be maximised if everything in this world is in harmony. \textbf{Peaceful,
symbiotic co-existence is the key to harmony in the world.}

Co-existence (\emph{saha-astitva}) is also the fourth constituent of the
\textbf{comprehensive human goal} --- co-existence in nature. In the context
of co-existence, `prosperity' means \textbf{holistic well-being and
happiness}, not material wealth accumulation.

\begin{keybox}
\textbf{\sffamily For the one-mark version (``give any two definitions'').}
Any two of the four numbered definitions above will do. The safest pair is
\emph{(1)} to exist together and in mutual tolerance, and \emph{(3)} a
relationship in which none of the parties is trying to destroy the other.
\end{keybox}

% ------------------------------------------------------------------------ 64
\Q{64}{Define the terms Gathansheel, Gathanpurna, Kriyapurnata and
  Acharanpurnata. Explain the holistic perception of harmony in
  existence.}{\tg{CIE 3 QB}\mmk{5}}

\ans
\begin{cmptable}{Term}{Meaning}
\textbf{Gathansheel} & The material units are \textbf{transformable}, and
  their composition keeps on changing --- hence these are
  \emph{gathansheel}. \\ \hline
\textbf{Gathanpurna} & The other category of units, the sentient `I', \textbf{do
  not transform} and are \textbf{complete in composition} --- hence
  \emph{gathanpurna}. \\ \hline
\textbf{Kriyapurnata} & \textbf{Completion of right understanding} in the human
  being. \\ \hline
\textbf{Acharanpurnata} & The \textbf{ability to live with complete
  understanding}. \\ \hline
\end{cmptable}

\textbf{Holistic perception of harmony in existence.} Existence is units in
space. Space is the empty area all around. The units are of two types:
\textbf{material (insentient)} and \textbf{conscious (the sentient `I')}. The
material units are changeful (with the activities of recognising and
fulfilment only), while the other kind of units are continuous (with the
activities of knowing, assuming, recognising and fulfilment).

The material units are available in two orders --- \textbf{material order} and
\textbf{pranic order}. In the material order, an atom combines with another
atom to form a molecule; a molecule similarly forms a molecular structure.
Molecular structures are found in two forms: \textbf{lumps and fluids}. Fluids
give nutrition to the pranic order. In the pranic order, the smallest units
are \textbf{plant cells}, which combine with other cells to form plants,
animal bodies and human bodies. The \textbf{co-existence of `I' with the
animal body becomes the animal order}, and the \textbf{co-existence of `I'
with the human body becomes the human order}.

On the material side, the transformation keeps taking place and it is
\textbf{cyclic} in nature. On the side of the sentient `I', the transitions are
\textbf{acyclic}. This implies that what we have understood \textbf{continues
to stay with us} --- we will never miss it. This is a transition in one
direction, and this is actually called \textbf{development (vikas)}.

So, existence is in the form of co-existence; it is in harmony. We don't have
to create this harmony --- it already exists; we only have to understand it in
order to be in it. This means that having the knowledge of the self (`I')
gives me the knowledge of \textbf{humane conduct} --- how to live in
existence, with the four orders. With this knowledge I can live with humane
conduct. \textbf{This is the pending task we have to complete.}

% ------------------------------------------------------------------------ 65
\Q{65}{Comment on the statement: ``Nature is limited and space is unlimited.''
  Differentiate between units and space. \emph{(Also set as:} Elaborate the
  terms `units' and `space' and mention their characteristics.\emph{)}}%
  {\tg{CIE 3 QB}\tg{Impr.\ CIE 2025}\tg{Impr.\ CIE 2026}\mmk{1--7}}

\ans Nature has four orders and there are units in each order. \textbf{Each
unit is limited in size.} The size ranges from being really small (an atom) to
really big (galaxies). Each and every unit is \textbf{finite and limited} in
size, be it the smallest particle or the biggest galaxy.

\textbf{Space, on the other hand, is unlimited.} Space has no `size'; unlike
units, it is not bounded. There is no beginning or end to space, as there is
to units.

\textbf{Illustration.} When we take a book, we know that it starts and
finishes --- the book is `limited' in size. When we take space, there is no
such thing. There is space behind us, inside us, between us and the book,
between the book and the earth, in the book, in every page of it, inside the
page, and beyond the earth --- all the way till we can imagine. Even if we say
space ends \emph{here} and there is ``nothing'' after that, that ``nothing''
is still space. We find that space \textbf{pervades}; it is all-pervading.
Units, on the other hand, are not all-pervading. This is precisely how we
recognise them as units.

\vspace{2pt}
\textbf{The five characteristics, side by side} --- this is the full form the
seven-mark version of the question wants:

\begin{cmptable}{Units (Nature)}{Space}
\textbf{Limited} (\emph{simit}) in size; bounded on six sides; countable and
  recognisable; not all-pervading &
\textbf{Unlimited} (\emph{asimit}); not bounded; no beginning or end;
  \textbf{all-pervading} (\emph{vyapaka}) \\ \hline
\textbf{Active} (\emph{kriyasil}) --- every unit of all four orders has
  activity, be it physical, physico-chemical or sentient. When something is
  active or has activity, we call it a \emph{unit}. &
\textbf{No-activity} (\emph{kriya-sunya}). There is empty space between you
  and this book and it has no activity. That is how we come to know of it.
  \\ \hline
\textbf{Energised} (\emph{urjit}) --- anything that is a unit has activity, and
  anything that has activity is energised. &
\textbf{Energy in equilibrium} (\emph{samya urja}) --- constant energy. This
  energy is available to all units. \\ \hline
\textbf{Each unit recognises and fulfils its relation with other units}
  (\emph{parasparata ko pahachanana, nirvaha karana}) --- the water recognises
  its relationship with the soil, the air with the soil, and so on. &
\textbf{Reflecting and transparent} (\emph{paradarsi}) --- every unit is
  reflected in the other units in space; because space is transparent there is
  no obstruction, so all units are able to recognise each other and are
  related being in space. \\ \hline
\textbf{Self-organised} (\emph{niyantrit}) --- every unit maintains its own
  organisation; no one organises it from outside. &
\textbf{Self-organisation is available} (\emph{niyantran upalabdha}) --- space
  is not itself self-organised; it is that in which self-organisation is
  available to units. \\ \hline
\end{cmptable}

\begin{keybox}
\textbf{\sffamily To sum up (the text's own summary).} There are two kinds of
realities in existence: \textbf{space, and units in space}. The units are in
co-existence with space and among themselves. Each unit is energised and
active in space, self-organised in space, and recognising and fulfilling its
relationship with other units in space. Space is in continuum: all-pervading
and no-activity. We are also units in space, and there is acceptance in us (in
`I') for self-organisation.
\end{keybox}

% ------------------------------------------------------------------------ 66
\Q{66}{Explain `energised' and `energy in equilibrium'.}{\tg{CIE 3 QB}\mmk{4}}

\ans What we normally call or consider as energy today is the \textbf{transfer
of energy}. For example, when we place water in a vessel on the stove, we say
the heat energy from the flame was transferred to the water in the vessel.

What about the water \emph{before} we put it on the stove --- was it
energised? What about the stove before we lit it? We may normally think ``no'',
but the fact is that it was. \textbf{Anything that is a unit has activity; and
anything that has activity is energised.} All the particles in the water and
the metal stove are active, very active, and energised. We can't see this
easily, or we don't see the `physical effects' of the unheated water or the
unlit metal stove, but they are still energised.

Space, on the other hand, is not a unit and has no activity. Hence we do not
say that space is energised; we say \textbf{``space is energy in
equilibrium''}, or that it is \textbf{constant energy}. All units are
energised in space, and this energy is available to all units.

\begin{keybox}
\textbf{Space is equilibrium energy; all units are in space; all units are
energised and active being in space.}
\end{keybox}

\section{D. Further Questions from the 2025--2026 Papers}

% ------------------------------------------------------------------------ 71
\Q{71}{Outline the 2 types of units with examples.}{\tg{Impr.\ CIE 2026}%
  \mmk{5}}

\ans Existence is units in space. Those units are categorised into
\textbf{two types}:

\begin{cmptable}{Material units (\emph{jada})}%
                {Conscious units (\emph{chaitanya})}
\textbf{Activities:} recognising and fulfilment \emph{only}. There is no
  assuming in them --- no potential to know or to assume. &
\textbf{Activities:} knowing, assuming, recognising and fulfilment. Currently
  assuming, recognising and fulfilment are predominant in humans; the capacity
  to \emph{know} is explored through natural acceptance. \\ \hline
\textbf{Temporary} in nature --- they undergo structural changes. Also called
  \textbf{gathansheel}. &
\textbf{Continuous} --- there is no structural change in them; there can only
  be a \emph{qualitative} change. Also called \textbf{gathanpurna}. \\ \hline
\textbf{Examples:} everything in the material order --- stones, minerals,
  soil, petrol; and everything in the plant\,/\,bio order --- plants, shrubs,
  grass; \emph{and also the animal body and the human body}. &
\textbf{Example:} the Self, `I' (\emph{Jivana}) --- in the animal order and in
  the human order. \\ \hline
\end{cmptable}

\textbf{Why the body counts as a material unit.} This is the point the
question is testing. If we consider the animal and human \emph{body}, we find
there is only recognition and fulfilment --- there is no assuming in the body.
The text's illustration: take a blade and cut your finger. If the blade is
sharper than your skin, the skin gets cut and bleeds --- no matter how many
times you try it, and no matter who does it. The skin does not respond
differently depending on who is cutting it.

But at the level of the Self, assumption \emph{is} involved. If you have gone
to a doctor's clinic, you know why the cut is being made and you cooperate. If
you are sitting in a bus and a stranger takes a knife to your skin, you are
alarmed and resist. The body behaves identically in both cases; only the
\emph{assumption} differs. Hence:

\begin{keybox}
At the level of the \textbf{body} there is only recognising and fulfilment.
\newline
At the level of \textbf{`I'} there is assuming and knowing, \emph{in addition
to} recognising and fulfilment.
\end{keybox}

\textbf{Why this classification is worth making.} If we know which units assume
and which do not, we know how to interact with them. Interacting with the
material order --- a stone --- we know no assumption is involved and can be
assured of its behaviour. Interacting with animals or humans, we know
assumption \emph{is} involved, and keeping this in consideration we can
improve our interaction and our relationship.

% ------------------------------------------------------------------------ 72
\Q{72}{Justify that the requirement of any order is already available in
  abundance in nature.}{\tg{Impr.\ CIE 2026}\mmk{5}}

\ans The claim has two halves, and both have to be established: that the
\emph{requirement} is limited, and that the \emph{availability} exceeds it.

\textbf{1. The requirement is limited.} A human being is the co-existence of
the Self (`I') and the Body. The needs of the Body --- physical facilities ---
are \textbf{quantitative and limited}: four chapattis, one bicycle. Beyond the
limit, consumption goes from \emph{necessary and tasteful} to
\emph{unnecessary but tasty} to \emph{unnecessary and tasteless} to
\emph{intolerable}. We have a common misconception today that our ``needs are
unlimited''; but on considering the needs of the body in order to maintain
health and right utilisation, it turns out to be \textbf{limited}.

\textbf{2. The availability is more than the requirement.} If you look around
in nature, the availability of facilities is \emph{more than our needs}. As
long as we are looking at the real needs of the Body, \textbf{there is enough
in nature already for all of us}. The text's own examples: there is more wheat
grown in the world than we can eat, more oxygen and air than we can breathe,
and more water than we can drink.

\textbf{3. This holds for the other orders too, by mutual fulfilment.} Each
order's requirement is met in abundance by the others, which is what
interconnectedness in nature means (\textbf{Q59}): the material order provides
nutrients to the pranic order in the form of soil and minerals, while the
pranic order decays and enriches the soil in return; the material order
provides the basis for the movement of all animals; water, oxygen and other
gases are available for both plants and animals; the plant order provides food
for animals, birds and fishes. The material order is itself the \emph{most
abundant} in nature. None of these orders denies the other, and none of them
runs short.

\textbf{4. Hence prosperity is possible.} It is \textbf{only} when we see that
the availability in nature is more than the limited needs of the body, and
that these can easily be fulfilled, that there emerges \textbf{a possibility
of prosperity} --- prosperity being the feeling of having \emph{more than
required} physical facilities (\textbf{Q15}).

\textbf{5. And hence the converse is the real problem.} Today we have assumed
our needs to be unlimited, by confusing `I' with the Body, and so we keep
running after ``unlimited amounts''. We fail to see that ``unlimited'' means
``having no end'', and so we try to accumulate endlessly --- a mad race. If we
can see the simple point that the needs of the Self and the needs of the Body
are different, and that the needs of the body are limited, there is the
possibility of prosperity for each one of us.

\begin{keybox}
\textbf{If the needs are assumed unlimited, then there is not enough even for
any ONE of us.} The shortage is manufactured by the assumption, not by nature.
\end{keybox}

% ------------------------------------------------------------------------ 73
\Q{73}{Summarize how the domain of consciousness helps in accomplishing
  development or permanent change.}{\tg{Impr.\ CIE 2026}\mmk{5}}

\ans \textbf{What development means.} Start with the question the text asks:
if you had `development', would you want it to be reversible or irreversible?
Would you want what you gained to stay, or to go back to where you started?
The answer is that we want development to be \textbf{irreversible, to stay
with us, to be permanent}. So if something we do is in a cycle --- if we go
back to where we started from --- \textbf{we would not call it development at
all}.

\textbf{Why the material domain cannot give it.} The material and pranic
orders are \textbf{cyclic by nature}. They undergo
composition\,/\,decomposition; the material units are \emph{gathansheel} ---
transformable, temporary, forever changing structure. No matter what we do
with the material and plant order, it will transform and go back. Hence
\textbf{we cannot `develop' there.} We shall keep losing that `development',
since it is cyclic. We can put things there to \emph{use} --- but use is not
development.

\textbf{Why the domain of consciousness can.} The Self (`I') is a
\textbf{conscious (chaitanya)} unit: it is \emph{gathanpurna} ---
\textbf{continuous}, with no structural change in it. The only change possible
in it is a \textbf{qualitative} change. On the material side the transformation
is cyclic; on the side of the sentient `I' the transition is \textbf{acyclic
--- in one direction only}. This implies that what we have understood
\textbf{continues to stay with us; we will never lose it}. That one-directional
qualitative transition is precisely what the text calls \textbf{development
(vikas)}.

\begin{keybox}
Change in `I' is \textbf{qualitative} and \textbf{permanent}. \textbf{This is
development.} If we want development, this is where we need to be working ---
for right understanding (knowledge) and right living.
\end{keybox}

\textbf{Where each domain should therefore be used.} The text sets out the
correct division of effort:

\begin{abul}
\item \textbf{For composition\,/\,construction, work on the material order} ---
      houses, implements, transport, communication; things we want to remain
      as constructed. Production here can stay within nature's framework of
      recycling.
\item \textbf{For growth, work on the plant\,/\,pranic order} --- we work to
      enrich nature, and there is enough in nature for all of us.
\item \textbf{For development, work on `I'} --- by knowing reality, knowing
      the entire existence, and living in accordance with it. Ensuring right
      understanding and right living in `I' ensures that our interaction with
      all four orders is in harmony, and that we are in harmony in ourselves.
\end{abul}

\textbf{The mistake we make today.} What we are doing in the name of
`development' is to focus on the material order and the plant order --- which
are cyclic --- so we are, in the text's words, \textbf{investing ourselves to
grow things that do not grow, and to develop things where it cycles back}.
Accumulation in the material order cannot be a substitute for the needs of
knowledge, understanding and relationship in `I'. Work on the material order
needs to be done \emph{in the light of} understanding in `I' and the needs of
`I'.

The completion of this development has two names in the text:
\textbf{kriyapurnata} --- the completion of right understanding --- and
\textbf{acharanpurnata} --- the ability to live with that complete
understanding (\textbf{Q64}).

% ------------------------------------------------------------------------ 74
\Q{74}{Explain the classification of units into Four Orders and the harmony
  among them with the help of a schematic diagram.}{\tg{Impr.\ CIE 2025}%
  \mmk{6}}

\ans \textbf{The classification.} Existence is nature submerged in space. All
the units of nature fall into two kinds --- material (\emph{jada}) and
conscious (\emph{chaitanya}) --- and these build up, in a self-organised
sequence, into the four orders:

\begin{abul}
\item In the \textbf{material order}, an atom combines with another atom to
      form a molecule; molecules form molecular structures, which exist either
      as \textbf{lumps (pinda)} or \textbf{fluids (rasa)}.
\item Fluids are the basis of the \textbf{plant cell}, and such cells combine
      to form plants --- this is the \textbf{pranic order}. The same cells go
      on to form the animal body and the human body.
\item The \textbf{animal order} is the co-existence of the animal body (pranic
      order) with the Self, `I'.
\item The \textbf{human order} is the co-existence of the human body (pranic
      order) with the Self, `I'.
\end{abul}

All these processes take place in a \textbf{self-organised, natural manner.
There is nothing controlling them}; they occur naturally in co-existence.

\vspace{6pt}
\begin{center}
\begin{tikzpicture}[
  font=\sffamily\scriptsize,
  ord/.style={draw=rulegrey,fill=panel,rounded corners=1pt,line width=0.7pt,
              align=center,inner sep=4pt,text width=27mm,minimum height=11mm},
  hum/.style={draw=ocre,fill=tint,rounded corners=1pt,line width=1pt,
              align=center,inner sep=4pt,text width=27mm,minimum height=11mm},
  lbl/.style={font=\sffamily\scriptsize\itshape,text=slatelight},
  mf/.style={{Stealth[length=4pt]}-{Stealth[length=4pt]},line width=0.7pt,
             draw=ocre},
  miss/.style={{Stealth[length=4pt]}-{Stealth[length=4pt]},line width=0.7pt,
               draw=slatelight,dashed}]

\node[ord] (m) at (0,0)    {\textbf{Material Order}\\ \emph{padartha avastha}\\
  soil, water, air, metals};
\node[ord] (p) at (4.1,0)  {\textbf{Pranic\,/\,Bio Order}\\ \emph{prana
  avastha}\\ plants, herbs, trees};
\node[ord] (a) at (8.2,0)  {\textbf{Animal Order}\\ \emph{jiva avastha}\\
  animals, birds};
\node[hum] (h) at (12.3,0) {\textbf{Human Order}\\ \emph{gyana avastha}\\
  human beings};

\draw[mf] (m) -- (p);
\draw[mf] (p) -- (a);
\draw[miss] (a) -- (h);

\draw[mf] (m.south) .. controls (2.05,-1.5) and (6.15,-1.5) .. (a.south);
\draw[miss] (p.south) .. controls (6.15,-2.3) and (10.25,-2.3) .. (h.south);
\draw[miss] (m.south) .. controls (4.1,-3.1) and (8.2,-3.1) .. (h.south);

\node[lbl] at (6.15,0.75) {mutual fulfilment --- ensured};
\node[lbl,text=ocre] at (10.25,0.75) {\textbf{?} not ensured by us};
\node[lbl] at (6.15,-3.45) {The three orders other than the human order are
  already in harmony and are fulfilling to the human order.};
\end{tikzpicture}
\end{center}

\textbf{The harmony among them.} The solid links above are relations of
\textbf{mutual fulfilment} that already hold; the dashed links are the ones
\textbf{we have failed to ensure}.

\begin{apts}
\item \textbf{Material $\leftrightarrow$ Pranic.} The material order provides
      nutrients as soil and minerals; the pranic order decays and enriches the
      soil in return, holds it against erosion, and produces
      oxygen\,/\,carbon dioxide.
\item \textbf{Material, Pranic $\leftrightarrow$ Animal.} The material order
      provides the basis for the movement of animals, birds and fishes; the
      animal order enriches the soil with its excreta and pollinates the
      flowers of the pranic order; the pranic order feeds the animals.
\item \textbf{The three orders $\rightarrow$ Human.} Each of these orders is
      fulfilling to the human order --- verifiable from the multiple uses we
      draw from them.
\item \textbf{Human $\rightarrow$ the three orders: missing.} We too have a
      natural acceptance to be mutually fulfilling to the other three, but
      \textbf{we are not able to ensure it}. We depend on the material order
      but pollute the soil and deplete the fossil fuels; we depend on plants
      but have deforested and destroyed species; we depend on animals but have
      made many extinct.
\end{apts}

\textbf{Conclusion.} There is interconnectedness and mutual fulfilment in all
the orders of nature \textbf{except the human order} --- and that single
missing link is the whole problem, and the whole task. We have not understood
the harmony that exists between the orders, nor our own needs properly, and
have consequently disturbed both ourselves and the balance among the other
three.
